\topsection[sec:threat]{Threat model}

The threat model links phase~1 (what must be protected) to phase~3 (what the design must enforce). The full
version, with entry points and the complete threat table, is \repofile{docs/04_threat_model.md}. The method is
an asset-centric analysis with a STRIDE classification. Each threat is mapped to the verifier check that stops
it and to the attack variant that demonstrates it.

\subsection{Trust assumption and trust boundaries}
One assumption carries the whole design. The signing keys (including the root key and the small pin file that
anchors the trust log), the policy and the verifier code are trusted. Everything else may be hostile
(Figure~\ref{fig:trust}). Under this assumption \sys\ detects tampering \emph{after} attestation. It does not
detect malice \emph{before} attestation.

\begin{figure}[!htb]
\centering
\begin{tikzpicture}[x=1cm,y=1cm,font=\sffamily\fontsize{7.9}{10.2}\selectfont,
  zone/.style={rounded corners=3pt,minimum width=4.75cm,minimum height=3.3cm,inner sep=0pt},
  ttl/.style={anchor=north west,font=\sffamily\bfseries\fontsize{8.8}{11}\selectfont},
  body/.style={anchor=north west,align=left,text width=4.3cm,text=cgink}]
\node[zone,fill=cgband] (u) at (0,0) {};
\node[zone,fill=cgband] (p) at (5.12,0) {};
\node[zone,draw=cgnavy,line width=0.9pt,fill=white] (t) at (10.24,0) {};
\node[ttl,text=cgnavy] at ($(u.north west)+(0.2,-0.15)$) {Hostile};
\node[body] at ($(u.north west)+(0.2,-0.68)$) {Data sources\\Outside model publishers and hubs\\Artifact storage and mirrors\\The network\\Manifests and log records at rest};
\node[ttl,text=cgnavy] at ($(p.north west)+(0.2,-0.15)$) {Not trusted for decisions};
\node[body] at ($(p.north west)+(0.2,-0.68)$) {Training and fine-tuning hosts\\The SQLite index (a convenience copy: never the source of a trust or revocation decision)\\Cached verification status};
\node[ttl,text=cgnavy] at ($(t.north west)+(0.2,-0.15)$) {Trusted computing base};
\node[body] at ($(t.north west)+(0.2,-0.68)$) {Private signing keys, one per role\\Root key and pin file of the trust log\\Policy file (who may sign what)\\Verifier and gateway code};
\end{tikzpicture}
\caption{Trust boundaries. The verifier recomputes every hash, digest and Merkle root instead of believing
stored values, and reads nothing from the index that it cannot confirm against signed evidence.}
\label{fig:trust}
\end{figure}

\subsection{Actors}
{\tablefont
\begin{xltabular}{\textwidth}{P{3.9cm}P{5.2cm}L}
\caption{Actors, their capabilities and their goals.}\label{tab:actors}\\
\toprule
\headrow\hd{Actor} & \hd{Capability} & \hd{Goal} \\
\midrule
\endfirsthead
\toprule
\headrow\hd{Actor} & \hd{Capability} & \hd{Goal} \\
\midrule
\endhead
\bottomrule
\endfoot
External attacker & Write access to artifact storage or a mirror & Modify or swap datasets and models \\
Malicious data contributor & Adds samples before attestation & Plant a backdoor that gets legitimately signed \\
Compromised training host & Runs code during training; may read the pipeline key & Alter weights, sign rogue models \\
Malicious model publisher & Distributes pretrained models & Ship a backdoored or code-executing model \\
Registry operator (insider) & Edits the database, the manifests or the log & Hide history, forge a status, undo a revocation \\
Legitimate signers & Hold the role keys & Not malicious, but can sign poisoned content in good faith \\
\end{xltabular}}

\subsection{Threats, controls and evidence}
Table~\ref{tab:threats} lists what the actors of Table~\ref{tab:actors} can do, with the control that addresses
each threat and the experiment variant that exercises it. The identifiers of the variants (A1, E5-C4 and so on) are those of Section~\ref{sec:eval}
and Appendix~\ref{app:variants}; the checks V0 to V9 are defined in Section~\ref{sec:design}.

{\tablefont\setlength{\tabcolsep}{4pt}
\begin{xltabular}{\textwidth}{P{4.25cm}P{0.9cm}P{3.9cm}P{1.25cm}P{1.75cm}L}
\caption{Threats, controls and evidence. STRIDE classes: S spoofing, T tampering, R repudiation, E elevation of privilege.}\label{tab:threats}\\
\toprule
\headrow\hd{Threat} & \hd{Class} & \hd{Control} & \hd{Check} & \hd{Shown by} & \hd{Detected} \\
\midrule
\endfirsthead
\toprule
\headrow\hd{Threat} & \hd{Class} & \hd{Control} & \hd{Check} & \hd{Shown by} & \hd{Detected} \\
\midrule
\endhead
\bottomrule
\endfoot
Dataset modified after signing & T & File hash and row-level Merkle root & V1, V2 & A1 to A4, E1 & yes \\
Poisoned data before signing & T & Content checks; accountability of the signer & V9 & E2, E3, E4 & in part \\
Model file substituted or altered & T, S & Hash bound in the signed payload & V1 & B1, B2 & yes \\
Model in a format that runs code & E & safetensors only; verify before load & V1, gateway & B3 & yes \\
Parent link removed or edited & T & Parent digests inside the signed payload & V3, V5 & C1, C3 & yes \\
History rewritten with a valid key & T & Hash links held by the children & V5 & C2 & yes \\
Model without dataset ancestry & T & Stage and parent policy & V8 & C4 & yes \\
Forged, corrupted or replayed signature & S & Ed25519 over the canonical payload & V3 & D1, D2, D6 & yes \\
Signature by an unknown key & S & Keys must be in the signed trust log & V3, V4 & D3, E7-B & yes \\
Key used outside its role & E & Role binding and policy & V4 & D5 & yes \\
Revoked key or artifact used & S & Signed revocation, propagated to descendants & V4, V6, V7 & D4, E5-A1, E5-A2, E5-B1 & yes \\
Index edited to undo a revocation or add a key & T & Signed trust log replayed at every verification & V4, V7 & E5-C1 to C5 & yes \\
Trust log truncated, edited or deleted & T & Root signatures, hash chain, pinned head; fail closed & V3, V4, V7 & E5-C6 to C8 & yes \\
File swapped between check and use & T & Read once, hash the bytes that are loaded & gateway & tests & yes \\
Backdoored model from a trusted publisher & T & Publisher key; evaluation & V3, V4 & E7-F & integrity yes, backdoor no \\
Private key stolen & S & Owner-only files, CI secrets & none & none & no (residual) \\
Training host lies in its provenance & T & Signed claims, role separation & none & none & no (residual) \\
Unauthenticated use of the API & S, E & None in this version & none & none & no (residual) \\
\end{xltabular}}

\subsection{Residual risks}
\begin{enumerate}
\item \textbf{Poison and backdoors present before attestation} pass every cryptographic check. This is the
  limit that Section~\ref{sec:eval} measures.
\item \textbf{Key compromise} yields perfectly valid signatures until the key is revoked. Mitigations are
  hardware-backed keys, short-lived certificates and fast revocation.
\item \textbf{A compromised builder} can lie in the provenance it signs. Without hardware attestation or
  reproducible training, nobody can prove which code produced the weights.
\item \textbf{The verifier and the pin file} are the trusted computing base. They must be distributed signed,
  and the pin protected like the root key.
\item \textbf{Freshness.} There is no trusted timestamp; the creation date of a manifest is a claim of its
  signer, and a verifier with a stale copy of the trust log would miss a recent revocation.
\item \textbf{The API has no authentication} in this version, and the demonstration endpoints that tamper on
  purpose must be disabled outside demonstrations.
\end{enumerate}
